% 出典: 2025ap_exam.pdf, 日本語版 PDF pp. 8--9.
\documentclass[11pt,a4paper]{article}
\usepackage[margin=23mm]{geometry}
\usepackage{fontspec}
\usepackage{xeCJK}
\setCJKmainfont{HaranoAjiMincho-Regular.otf}
\setCJKsansfont{HaranoAjiGothic-Medium.otf}
\usepackage{amsmath,amssymb}
\usepackage{tikz}
\setlength{\parindent}{1em}
\setlength{\parskip}{0.35em}
\title{物理工学専攻 院試解答集\\2025年 物理学 第4問}
\author{政岡凜太郎}
\date{}
\begin{document}
\maketitle

\section*{[1.1] 電子の運動方程式}
電子の電荷は $-e$ なので、電場による力は $-e\boldsymbol E$ である。ポテンシャルの勾配を加えると、減衰のない運動方程式は
\[
\boxed{m\ddot{\boldsymbol u}=-\boldsymbol\nabla_{\boldsymbol u}V-e\boldsymbol E
=-m\begin{pmatrix}\omega_\alpha^2u_x\\\omega_\alpha^2u_y\\\omega_\beta^2u_z\end{pmatrix}
-e\boldsymbol E_0e^{-i\omega t}.}
\]
成分ごとには $m\ddot u_j+m\omega_j^2u_j=-eE_j$ であり、$\omega_x=\omega_y=\omega_\alpha$, $\omega_z=\omega_\beta$ と置ける。

\section*{[1.2] 誘電率テンソル}
定常解を $u_j(t)=u_{j,0}e^{-i\omega t}$ と置く。$\ddot u_j=-\omega^2u_j$ だから [1.1] より
\[
u_j=-\frac{e}{m(\omega_j^2-\omega^2)}E_j.
\]
共振点 $\omega_j=\omega$ は問題の条件で除かれている。$\boldsymbol P=-\rho e\boldsymbol u$ を $\boldsymbol D=\epsilon_0\boldsymbol E+\boldsymbol P$ に代入すると
\[
\boxed{\tilde\epsilon=
\begin{pmatrix}
\displaystyle\epsilon_0+\frac{\rho e^2}{m(\omega_\alpha^2-\omega^2)}&0&0\\
0&\displaystyle\epsilon_0+\frac{\rho e^2}{m(\omega_\alpha^2-\omega^2)}&0\\
0&0&\displaystyle\epsilon_0+\frac{\rho e^2}{m(\omega_\beta^2-\omega^2)}
\end{pmatrix}.}
\]
電子密度を $0$ にすると真空の $\epsilon_0\mathbf I$ に戻る。分母の符号は共振周波数との大小に依存し、[2] の $\epsilon_n>\epsilon_a>\epsilon_0$ は別途与えられた条件である。

\section*{[2] 座標と二つの透過波}
原本図1に合わせ、境界を $z=0$、結晶を $z>0$ とする。$x$ は右、$z$ は図で下向き、$y$ は紙面手前を向く。以下の角は $z$ 軸から右側へ測り、透過波の $k_z$ は正とする。
\begin{center}
\begin{tikzpicture}[x=1cm,y=1cm,>=stealth,scale=0.95]
  \draw[gray,thick] (-3.5,0)--(3.7,0);
  \draw[->] (0,2.5)--(0,-2.5) node[below] {$z$};
  \draw[->] (0,0)--(3.6,0) node[right] {$x$};
  \node[left] at (-2.5,0.38) {真空 $z<0$};
  \node[left] at (-2.5,-0.38) {結晶 $z>0$};
  \draw (0.62,2.0) circle (0.12);
  \fill (0.62,2.0) circle (0.035);
  \node[right] at (0.80,2.0) {$y$：紙面手前};
  \draw[->,thick] (-2.1,1.9)--(0,0) node[pos=0.44,above left] {$\boldsymbol k_i$};
  \draw[->,thick] (0,0)--(1.8,1.65) node[near end,above] {$\boldsymbol k_r$};
  \draw[->,thick,blue] (0,0)--(0.70,-2.15) node[below] {$\boldsymbol k_{t1}$（S）};
  \draw[->,thick,red] (0,0)--(1.58,-2.15) node[right] {$\boldsymbol k_{t2}$（P）};
  \draw (0,-0.77) arc[start angle=-90,end angle=-72,radius=0.77];
  \node at (0.52,-0.90) {$\phi_{t1}$};
  \draw (0,-1.12) arc[start angle=-90,end angle=-54,radius=1.12];
  \node at (1.18,-1.15) {$\phi_{t2}$};
\end{tikzpicture}
\end{center}
この図は進行方向の関係を示す模式図で、矢印の長さは波数の大きさを表さない。

\section*{[2.1] 境界に沿う位相の一致}
境界 $z=0$ では各平面波の位相は $k_xx-\omega t$ となる。接線方向の電場・磁場の連続条件がすべての $x,t$ で成り立つには、非零の入射・反射・透過振幅の位相が一致しなければならない。したがって
\[
k_{i,x}=k_{r,x}=k_{t1,x}=k_{t2,x}.
\]
図の斜入射の透過枝では $k_{tj,x}=|\boldsymbol k_{tj}|\sin\phi_{tj}>0$ であり、垂直入射ではこの成分は零となる。位相一致から
\[
\boxed{\sin\phi_{t1}=\frac{|\boldsymbol k_i|}{|\boldsymbol k_{t1}|}\sin\phi_i,
\qquad
\sin\phi_{t2}=\frac{|\boldsymbol k_i|}{|\boldsymbol k_{t2}|}\sin\phi_i.}
\]

\section*{[2.2] 平面波の電場と電束密度}
$e^{i(\boldsymbol k\cdot\boldsymbol r-\omega t)}$ の共通因子を外すと、問題の Maxwell 方程式は
\[
\boldsymbol k\times\boldsymbol E=\mu_0\omega\boldsymbol H,
\qquad
\boldsymbol k\times\boldsymbol H=-\omega\boldsymbol D
\]
となる。第一式を第二式に代入し、三重積を展開すれば
\begin{align*}
\boldsymbol k\times(\boldsymbol k\times\boldsymbol E)
&=-\mu_0\omega^2\boldsymbol D\\
&=(\boldsymbol k\cdot\boldsymbol E)\boldsymbol k-|\boldsymbol k|^2\boldsymbol E.
\end{align*}
ゆえに $\boldsymbol e=\boldsymbol k/|\boldsymbol k|$ を用いて
\[
\boxed{\boldsymbol D=\frac{|\boldsymbol k|^2}{\mu_0\omega^2}
\left[\boldsymbol E-(\boldsymbol e\cdot\boldsymbol E)\boldsymbol e\right].}
\]
この式から $\boldsymbol k\cdot\boldsymbol D=0$ も直ちに分かる。異方性媒質では一般に $\boldsymbol k\cdot\boldsymbol E=0$ とは限らない。

\section*{[2.3] S偏光の波数}
S偏光では $\boldsymbol E=E_y\boldsymbol e_y$ が入射面に垂直なので、$\boldsymbol e\cdot\boldsymbol E=0$。一方、結晶中では $D_y=\epsilon_nE_y$ である。[2.2] の $y$ 成分を非零の $E_y$ で割ると
\[
\epsilon_n=\frac{|\boldsymbol k_{t1}|^2}{\mu_0\omega^2},
\qquad
\boxed{|\boldsymbol k_{t1}|=\omega\sqrt{\mu_0\epsilon_n}.}
\]

\section*{[2.4] S偏光の光線方向}
Maxwell の式から $\boldsymbol H_{t1,0}=(\boldsymbol k_{t1}\times\boldsymbol E_{t1,0})/(\mu_0\omega)$。$\boldsymbol k_{t1}\cdot\boldsymbol E_{t1,0}=0$ なので三重積により
\[
\boldsymbol E_{t1,0}\times\boldsymbol H_{t1,0}
=\frac{|\boldsymbol E_{t1,0}|^2}{\mu_0\omega}\boldsymbol k_{t1}.
\]
係数は正であり、光線は波数と平行である。したがって求める角は
\[
\boxed{\phi_{t1}.}
\]

\section*{[2.5] P偏光の波数}
$\boldsymbol k_{t2}=k(\sin\phi_{t2},0,\cos\phi_{t2})$、$\boldsymbol E_{t2,0}=(E_x,0,E_z)$ と置く。$\boldsymbol k\cdot\boldsymbol D=0$ と $D_x=\epsilon_nE_x$, $D_z=\epsilon_aE_z$ より
\[
\epsilon_n\sin\phi_{t2}E_x+\epsilon_a\cos\phi_{t2}E_z=0,
\qquad
\frac{E_z}{E_x}=-\frac{\epsilon_n}{\epsilon_a}\tan\phi_{t2}.
\]
ここで $E_x\ne0$ であり、$\cos\phi_{t2}>0$ の透過枝を使った。[2.2] を $x$ 成分に適用すると
\[
\epsilon_nE_x=\frac{k^2}{\mu_0\omega^2}
\left(\cos^2\phi_{t2}E_x-\sin\phi_{t2}\cos\phi_{t2}E_z\right).
\]
上の比を代入し、正の平方根を採用して
\[
\boxed{|\boldsymbol k_{t2}|=\omega\sqrt{\frac{\mu_0\epsilon_n\epsilon_a}
{\epsilon_a\cos^2\phi_{t2}+\epsilon_n\sin^2\phi_{t2}}}.}
\]
$\phi_{t2}=0$ なら電場は $x$ 方向で $k=\omega\sqrt{\mu_0\epsilon_n}$ となり、[2.3] に一致する。

\section*{[2.6] P偏光の光線方向}
$\boldsymbol H_{t2,0}=H_y\boldsymbol e_y$ であり、Maxwell の式から
\[
H_y=\frac{k}{\mu_0\omega}
\left(\cos\phi_{t2}E_x-\sin\phi_{t2}E_z\right).
\]
$E_z/E_x\le0$ かつ $\cos\phi_{t2}>0$ なので $H_y/E_x>0$。したがって
\[
\boldsymbol E_{t2,0}\times\boldsymbol H_{t2,0}
=H_y(-E_z,0,E_x),
\qquad
\tan\theta_{t2}=\frac{-E_zH_y}{E_xH_y}.
\]
[2.5] の電場成分比を使えば
\[
\boxed{\tan\theta_{t2}=\frac{\epsilon_n}{\epsilon_a}\tan\phi_{t2}.}
\]
$\epsilon_n>\epsilon_a$ かつ $\phi_{t2}>0$ なら光線は位相法線より $x$ 側へ傾く。$\epsilon_n=\epsilon_a$ の等方極限では二方向は一致する。

\end{document}
